\documentclass[11pt]{article}

\usepackage[letterpaper,margin=1in]{geometry}
\usepackage{amsmath,amssymb}
\usepackage{booktabs}
\usepackage{longtable}
\usepackage{array}
\usepackage{xcolor}
\usepackage{listings}
\usepackage{url}
\usepackage{hyperref}
\usepackage{titlesec}
\usepackage{parskip}
\usepackage{fancyhdr}
\usepackage{enumitem}

\hypersetup{
  colorlinks=true,
  linkcolor=black,
  urlcolor=blue!50!black,
  citecolor=black,
  pdftitle={MAC Array --- Design Specification},
  pdfauthor={LonghornSilicon},
  pdfsubject={MAC array reference model, version 0.1}
}

\titleformat{\section}{\normalfont\Large\bfseries}{\thesection.}{0.6em}{}
\titleformat{\subsection}{\normalfont\large\bfseries}{\thesubsection}{0.6em}{}
\titleformat{\subsubsection}{\normalfont\normalsize\bfseries}{\thesubsubsection}{0.6em}{}

\pagestyle{fancy}
\fancyhf{}
\lhead{\small MAC Array Design Specification}
\rhead{\small \texttt{mac-array-v0.1}}
\cfoot{\thepage}
\renewcommand{\headrulewidth}{0.4pt}

\definecolor{codebg}{gray}{0.97}
\definecolor{codekey}{rgb}{0.13,0.29,0.53}
\definecolor{codecom}{rgb}{0.25,0.5,0.25}
\lstset{
  basicstyle=\ttfamily\small,
  backgroundcolor=\color{codebg},
  frame=single,
  framerule=0.4pt,
  rulecolor=\color{black!30},
  xleftmargin=0.5em,
  xrightmargin=0.5em,
  aboveskip=0.6em,
  belowskip=0.6em,
  showstringspaces=false,
  breaklines=true,
  keywordstyle=\color{codekey}\bfseries,
  commentstyle=\color{codecom}\itshape,
  numbers=none
}

\setlist{itemsep=2pt,topsep=3pt}

\begin{document}

%==============================================================================
% Title page
%==============================================================================
\thispagestyle{empty}
\begin{center}
\vspace*{1.2in}
{\Huge\bfseries MAC Array}\\[0.4em]
{\LARGE\bfseries Design Specification}\\[1.2em]

{\large \textsc{LonghornSilicon LLM Inference Accelerator}}\\[0.4em]
{\large Attention Compute Unit (ACU) --- MAC array sub-block}\\[0.3em]
{\normalsize (Block 1 of 4 at the chip level; second of ACU's two sub-blocks)}\\[2em]

\rule{0.7\textwidth}{0.4pt}\\[1em]

\begin{tabular}{rl}
\textbf{Version}:    & \texttt{mac-array-v0.1} (first public draft) \\
\textbf{Date}:       & 2026-05-14 \\
\textbf{Status}:     & Reference model complete; RTL pending \\
\textbf{Author}:     & Chaithu Talasila, The University of Texas at Austin \\
\textbf{Reference}:  & \texttt{LonghornSilicon/attention-compute-unit} \\
\end{tabular}\\[1em]
\rule{0.7\textwidth}{0.4pt}

\vfill
\begin{minipage}{0.85\textwidth}
\small
\textbf{Scope.} This document describes the externally-visible behavior of
the MAC array block of the LonghornSilicon ACU. Decisions marked
\textbf{frozen} will not be revisited in v0.2; decisions marked
\textbf{open} may change as the RTL lands and the compiler team's
requirements clarify. The accompanying bit-accurate C/C++/Python
reference implementations are the executable form of this specification.
\end{minipage}

\vspace{0.3in}
\end{center}

\clearpage

%==============================================================================
\tableofcontents
\clearpage

%==============================================================================
% 1. Overview
%==============================================================================
\section{Overview}
\label{sec:overview}

The MAC array is the compute engine of the ACU. After the precision
controller (see \texttt{pc-isa-0.1}) decides INT8 or FP16 for a tile,
the MAC array executes the chosen matmul. The compiler emits
\texttt{matmul\_int8} or \texttt{matmul\_fp16} per tile based on the
gate's decision.

\medskip
\noindent\textbf{Functional summary.} For row-major inputs
$A \in \mathbb{R}^{M \times K}$, $B \in \mathbb{R}^{K \times N}$, the
block computes $C = A \cdot B \in \mathbb{R}^{M \times N}$ at either
INT8 or FP16 precision. Output precision matches input precision in
the FP16 path; output precision is widened to int32 in the INT8 path
(the compiler is responsible for any downcast or requantization).

%==============================================================================
% 2. Operations exposed
%==============================================================================
\section{Operations Exposed}
\label{sec:ops}

The block exposes four logical operations to the compiler:

\begin{table}[ht]
\centering
\small
\caption{MAC array operations.}
\label{tab:ops}
\begin{tabular}{@{}l p{3.4cm} l p{5cm}@{}}
\toprule
Operation                & Inputs                & Output                & Use \\
\midrule
\texttt{matmul\_int8}    & int8 $[M \!\times\! K]$, int8 $[K \!\times\! N]$ & int32 $[M \!\times\! N]$ & INT8 attention / MLP path \\
\texttt{matmul\_fp16}    & fp16 $[M \!\times\! K]$, fp16 $[K \!\times\! N]$ & fp16 $[M \!\times\! N]$ & FP16 attention / MLP path \\
\texttt{query}           & ---                   & \texttt{MacArrayInfo} & Query precisions, accumulator widths, PE grid \\
\texttt{estimate}        & $M, K, N$, dtype      & cycles + energy (pJ)  & Cost model for compiler-side scheduling \\
\bottomrule
\end{tabular}
\end{table}

\noindent
The compiler emits one of \texttt{matmul\_int8} or \texttt{matmul\_fp16}
for every attention sub-tile, dispatched on the precision controller's
decision. See the worked lowering example in
\path{docs/isa/precision_controller_isa.pdf} \S4.1 for the full
attention-tile flow.

%==============================================================================
% 3. Data type semantics
%==============================================================================
\section{Data Type Semantics}
\label{sec:dtypes}

\subsection{INT8 Path (Frozen)}
\label{sec:int8}

\begin{itemize}
\item Storage: signed two's-complement 8-bit (\texttt{int8\_t},
      range $-128 \ldots +127$).
\item Multiplication: signed 8-bit $\times$ signed 8-bit $\rightarrow$
      exact 16-bit product.
\item Accumulation: 32-bit signed accumulator (\texttt{int32\_t}).
      For $K$ up to $65\,536$ inputs the accumulator cannot overflow
      (worst-case $K \cdot 127 \cdot 127 \approx 1.06 \cdot 10^9$,
      fits in 31 bits).
\item No saturation, no requantization at the accumulator boundary;
      the compiler is responsible for any downcast / requantization
      step after the matmul.
\end{itemize}

\subsection{FP16 Path (Open --- v0.1 Simplification)}
\label{sec:fp16}

\begin{itemize}
\item Storage: IEEE-754 binary16 (\texttt{half}, 1 sign + 5 exponent +
      10 mantissa).
\item Multiplication: half $\times$ half $\rightarrow$ round-to-nearest-
      even $\rightarrow$ half.
\item Accumulation: fp32 accumulator (matches typical attention
      hardware).
\item Final output: rounded back to fp16 (round-to-nearest-even).
\item \textbf{v0.1 caveat}: the reference model uses
      \texttt{float} internally and rounds back to fp16 between
      operations. The eventual hardware may use a denser accumulator
      format (e.g., bfloat16) or skip the intermediate rounding. The
      compiler should not depend on bit-exact parity with v0.1's FP16
      path until v0.2 freezes the choice.
\end{itemize}

\subsection{Why Both Paths (Frozen)}
\label{sec:bothpaths}

The precision controller decides INT8 or FP16 per tile. The MAC array
therefore needs to dispatch to either path on a tile-by-tile basis. A
unified mixed-precision path is out of scope for v0.1 because the
cost and energy savings of INT8 only materialize when the entire
matmul stays in INT8.

%==============================================================================
% 4. Shape support
%==============================================================================
\section{Shape Support (Frozen)}
\label{sec:shapes}

\begin{itemize}
\item Arbitrary $M \times K \times N$ (no power-of-two requirement).
\item Recommended tile shape: $64 \times 64 \times 64$ (matches the
      precision controller's default tile).
\item A typical FlashAttention-style block computes:
      \begin{itemize}
      \item $Q \cdot K^{\top}$: $64 \times 64 \times 64$
      \item $\text{score} \cdot V$: $64 \times 64 \times 64$
      \end{itemize}
\item The reference is shape-agnostic --- the compiler picks the
      shape, the reference produces the correct output.
\end{itemize}

%==============================================================================
% 5. Concurrency
%==============================================================================
\section{Concurrency and Streaming (Open)}
\label{sec:concurrency}

The v0.1 reference is synchronous: \texttt{matmul\_int8(...)} blocks
until the result is in $C$. The eventual hardware will likely expose:

\begin{enumerate}
\item An asynchronous ``issue'' call that returns a token;
\item A ``wait'' call on the token;
\item Double-buffered input/output regions so the compiler can pipeline.
\end{enumerate}

For v0.1 the compiler can model these by wrapping the synchronous
calls in its own scheduler. We will revisit when the MAC array RTL
lands.

%==============================================================================
% 6. Cost model
%==============================================================================
\section{Cost Model (v0.1 Placeholders)}
\label{sec:cost}

\texttt{MacArrayInfo} exposes:

\begin{table}[ht]
\centering
\small
\caption{Cost-model parameters baked into v0.1.}
\label{tab:cost}
\begin{tabular}{@{}l l p{8cm}@{}}
\toprule
Parameter                & Value (v0.1) & Notes \\
\midrule
\texttt{pe\_grid\_m}     & 8            & PE grid rows (chip: 8$\times$8 = 64 PEs) \\
\texttt{pe\_grid\_n}     & 8            & PE grid cols (chip: 8$\times$8 = 64 PEs) \\
\texttt{int8\_throughput}& 64 MACs/cyc  & One INT8 MAC per PE; 64 PEs $\to$ 128 GOPS at 1\,GHz \\
\texttt{fp16\_throughput}& 16 MACs/cyc  & FP16 PE is $\sim$$4\times$ INT8 area, so $\frac{1}{4}$ the count \\
\texttt{pipeline\_depth} & 16 cycles    & Fixed startup latency \\
\bottomrule
\end{tabular}
\end{table}

\noindent
\texttt{estimate(M, K, N, dtype)} returns:
\begin{align*}
\text{cycles}     & \approx \left\lceil \frac{M \cdot N \cdot K}{\text{throughput}} \right\rceil + \text{pipeline\_depth} \\
\text{energy}_{pJ} & \approx M \cdot N \cdot K \cdot \text{per\_op\_energy}_{pJ}
\end{align*}

\noindent
Per-op energy values: $0.5$ pJ for INT8 MACs and $2.5$ pJ for FP16
MACs (rough numbers from published 16FFC accelerator data).

These are \textbf{placeholders} for the compiler's scheduler. The
numbers will tighten once the MAC array RTL is written and synthesized;
today they are a reasonable order-of-magnitude estimate. The compiler
should use them for relative comparison (INT8 vs FP16, which shape is
faster) rather than absolute timing.

%==============================================================================
% 7. What's NOT in v0.1
%==============================================================================
\section{What Is Not in v0.1}
\label{sec:exclusions}

\begin{itemize}
\item Sparsity acceleration (skip-zero PEs).
\item Mixed-precision MAC (e.g., INT4 $\times$ INT8 within a single op).
\item Batched matmul as a single op (compiler can loop).
\item Hardware fault injection / ECC.
\item Power-state management.
\end{itemize}

These are all out of scope for the reference model until the chip
architecture commits to them.

%==============================================================================
% 8. How the compiler binds this
%==============================================================================
\section{Compiler Binding}
\label{sec:binding}

Same template as the precision controller. Three layers:

\begin{enumerate}
\item \textbf{C++ class} (\texttt{lhsi::mac::MacArray}) for MLIR / TVM
      C++ backends.
\item \textbf{extern "C" API} (\texttt{lhsi\_mac\_*}) for plain-C
      runtimes and FFI.
\item \textbf{Python reference} (\texttt{mac\_array\_ref.py}) for
      compiler-side prototyping and unit tests.
\end{enumerate}

All three pass identical test vectors. If the C++ and Python disagree,
that is a bug --- open an issue.

%==============================================================================
% 9. Open questions
%==============================================================================
\section{Open Questions}
\label{sec:open}

Bring these to the next compiler-team sync; the answers shape v0.2.

\begin{enumerate}
\item Does your compiler emit individual matmuls, or does it expect a
      fused \texttt{attention(Q, K, V)} op? Affects whether we should
      add a fused-attention reference operation.
\item Do you need cycle-level cost estimates, or is operation-count
      estimation sufficient for your scheduler?
\item For FP16: do you expect bit-exact parity with our hardware's
      FP16 rounding, or just numerically-close output? Affects v0.2's
      rounding-mode decision.
\item Async issue + token model --- do you want this in v0.1 (we
      would add it now) or v0.3 (after RTL exists)?
\end{enumerate}

%==============================================================================
% 10. Verification status
%==============================================================================
\section{Verification Status}
\label{sec:verif}

The v0.1 reference passes the following checks
(see \path{sw/reference_model/test_mac_array_ref.cpp} and the Python
equivalent).

\begin{table}[ht]
\centering
\small
\caption{MAC array v0.1 verification matrix.}
\label{tab:verif}
\begin{tabular}{@{}l l@{}}
\toprule
Test                                            & Status \\
\midrule
INT8 matmul vs.\ naive int64 reference (random shapes) & exact match \\
FP16 matmul vs.\ naive double reference (random shapes) & within $\pm 5 \cdot 10^{-3}$ relative \\
Edge cases: zero, identity, signed extremes      & pass \\
extern "C" API matches C++ class                 & pass \\
fp16 round-trip helpers (float $\leftrightarrow$ bits) & pass \\
Cost-estimate sanity (FP16 slower / costlier than INT8) & pass \\
\bottomrule
\end{tabular}
\end{table}

%==============================================================================
% 11. Change log
%==============================================================================
\section{Change Log}
\label{sec:changelog}

\begin{itemize}
\item \texttt{mac-array-v0.1} (2026-05-14): First public draft.
      Reference model complete in Python, C++ class, and extern "C"
      shim. Verification matrix in \S\ref{sec:verif} passes. RTL
      pending.
\end{itemize}

\vspace{0.5in}
\noindent\rule{\textwidth}{0.4pt}\\
\noindent\small\textit{LonghornSilicon ACU MAC Array Design Specification.
This document is open and may be redistributed. The accompanying
reference implementations are at
\path{sw/reference_model/mac_array_ref.{hpp,cpp,py}} in
\texttt{LonghornSilicon/attention-compute-unit}.}

\end{document}
